% THE NEXT ROW IS LOOKED FOR IN THE ALIGNMENT'S OWN MODE.
%
% An alignment reads its rows one at a time, and before each it looks for
% what comes next -- expanding as it looks. That looking happens in the
% alignment's own list, not in the mode the row before it ran in: internal
% vertical for \halign, restricted horizontal for \valign. So an \iftrue
% standing at the head of an entry is traced there, and the entry's own
% mode only begins with the first thing the entry really puts down.
%
%   \halign{#&#\cr\iftrue A\fi&B\cr}
%
%   {internal vertical mode: \iftrue}       the peek, in the alignment's list
%   {true}
%   {restricted horizontal mode: the letter A}   the entry proper
%
% HSTeX stayed in whatever mode the last entry had left, so the peek drew
% no mode at all.
%
% What stood before the rows stands again once they are done, which is
% what the last case here reads.
\catcode`\{=1 \catcode`\}=2 \catcode`\#=6 \catcode`\&=4
\tracingonline=1 \hsize=100pt \parindent=0pt \tabskip=0pt
\hbadness=10000 \vbadness=10000 \hfuzz=1000pt \vfuzz=1000pt
\message{[A an expandable at the head of a halign entry]}
\tracingcommands=2
\halign{#&#\cr\iftrue A\fi&B\cr}
\tracingcommands=0
\message{[B and at the head of the second row]}
\tracingcommands=2
\halign{#&#\cr A&B\cr\iftrue C\fi&D\cr}
\tracingcommands=0
\message{[C the same for valign]}
\tracingcommands=2
\noindent A\valign{#\cr\iftrue B\fi\cr}C\par
\tracingcommands=0
\message{[D what stands after the rows are done]}
\tracingcommands=2
\halign{#&#\cr A&B\cr} \relax
\tracingcommands=0
\message{[done]}
\end
